\clearpage
\thispagestyle{empty}
\pdfbookmark[0]{Magnetisches Observatorium in Göttingen. August 1834. Ausgabe und Quellen}{observatory-1834-edition}
\begin{center}\Large\textbf{Magnetisches Observatorium in Göttingen}\end{center}

Der vollständige Bericht vom August 1834 über das Göttinger magnetische Observatorium mit beiden Tabellen ist eine KI-Übertragung des deutschen Textes in Carl Friedrich Gauß’ \textit{Werken}, Band V (Göttingen, 1867), S.~519--525, in modernes Deutsch. Der Satz wurde mit KI erstellt. Der Bericht beschreibt das Gebäude und die Instrumente, die verabredeten Beobachtungen der magnetischen Deklination, die Messungen der Intensität des Erdmagnetismus und Wilhelm Webers Drahtverbindung für galvanische Versuche und telegrafische Signale.

Der Bericht erschien in den \textit{Göttingischen gelehrten Anzeigen}, 128.~Stück, 9. August 1834, S.~1265--1274. Eine gekürzte Fassung erschien im \textit{Polytechnischen Journal}, Band~55 (1835), Miszellen, ab S.~393. Einen gekürzten englischen Auszug zur Telegrafie enthält der ohne Verfassernamen erschienene Beitrag \textit{Sketch of William Edward Weber}, \textit{Popular Science Monthly}, Band~40, Februar 1892.

\medskip\noindent
\href{https://gdz.sub.uni-goettingen.de/id/PPN236006339}{Deutsche Werkausgabe (1867)}\par
\noindent\href{https://doi.org/10.1007/978-3-642-49319-5_28}{Deutscher Nachdruck (1877)}\par
\noindent\href{https://dingler.bbaw.de/articles/mi055069_1.html}{Gekürzte deutsche Fassung (1835)}\par
\noindent\href{https://en.wikisource.org/w/index.php?title=Popular_Science_Monthly/Volume_40/February_1892/Sketch_of_William_Edward_Weber\&oldid=8847389}{Englischer Auszug (1892)}
\clearpage
